\section{Two-day literature and hypothesis-freeze protocol}
\label{app:twoday}

The initial exploration is limited to two days.  Its purpose is not to survey the
entire electron--phonon literature, but to decide whether the channel-selective
hypothesis is physically defensible and to freeze the first decisive comparison.

\subsection*{Day 1: establish constraints and collect contrasts}

\begin{enumerate}
  \item \textbf{Morning: foundations.}  Verify the DFPT response convention,
  physical quasiparticle vertex, proper versus screened vertex, Ward identity, and
  static/dynamic limit ordering from primary sources
  \citep{Baroni2001,Giustino2017,Hedin1965,Nambu1960}.
  \item \textbf{Midday: UEG.}  Reproduce the algebra of
  \cref{eq:matchingcondition,eq:propervertexmatching}; verify that the complete
  $K_{\mathrm{total}}\simeq1$ rather than applying the proper-vertex ratio twice;
  extract the reported $q$, $r_s$, temperature, and
  error ranges from \citet{Cai2025}; list which statements are numerical and which
  are exact.
  \item \textbf{Afternoon: correlated contrasts.}  Extract mode definitions,
  eigenvectors, parameters, normalization, static vertices, and frequency
  dependence for SrVO$_3$ and CaCuO$_2$.

  The primary source is \citet{Abramovitch2025}.  Add CoO and one GWPT case to
  identify failure mechanisms not represented by local DMFT.
  \item \textbf{Evening gate.}  Produce a one-page evidence table with at least
  one supporting and one challenging case for each proposed channel.
\end{enumerate}

\subsection*{Day 2: turn the idea into a falsifiable first experiment}

\begin{enumerate}
  \item \textbf{Morning: define observables.}  Freeze the localized subspace,
  operator normalization, Fermi-window weights, target $g$ convention, and basis
  robustness checks.
  \item \textbf{Midday: compare alternatives.}  Test at least three hypotheses on
  the evidence table: correlation strength alone, quasiparticle residue alone,
  and channel weight times channel susceptibility.  Record which observations
  each cannot explain.
  \item \textbf{Afternoon: design the first calculation.}  Select a SrVO$_3$
  Jahn--Teller/breathing contrast and a UEG density control.  Represent the
  finite-R mode in the phase-correct $2\times2\times2$ real supercell or equivalent
  standing-wave basis before assigning channel weights.  Specify input data, code
  path, acceptance thresholds, and the result that would reject the channel hypothesis.
  \item \textbf{Final gate.}  Freeze one paragraph stating the physical
  hypothesis, one equation for the first correction model, one benchmark table,
  and one stop/go decision for implementation.
\end{enumerate}

\subsection*{Required outputs after 48 hours}

\begin{itemize}
  \item a source-checked evidence matrix with no mixed normalization conventions;
  \item a claim ledger labeling exact results, established methods, numerical
  evidence, and hypotheses;
  \item the first two material/mode comparisons and their input availability;
  \item a minimal operator-decomposition prototype specification; and
  \item an explicit decision: implement, revise the hypothesis, or stop.
\end{itemize}

The literature sprint is successful only if it reduces uncertainty about the
physical hypothesis.  The number of papers read is not itself a success metric.
